\section{Programação Dinâmica Dual Estocástica}
\label{sec:pdde_literatura}

A Programação Dinâmica Dual Estocástica (PDDE), conhecida na literatura internacional por \emph{Stochastic Dual Dynamic Programming} (SDDP), é um método de decomposição para problemas de Programação Estocástica Multiestágio introduzido por \citeonline{pereira1991}. O método foi desenvolvido originalmente para o planejamento de operação de sistemas hidrotérmicos e, desde então, tem sido estendido a outros domínios com estrutura multiestágio \cite{dowson2021}.

\subsection{Aproximação da Função de Valor por Cortes de Benders}
\label{subsec:cortes_benders}

A PDDE evita a construção explícita da árvore de cenários ao substituir a função de valor futuro $V_{t+1}(\cdot)$, desconhecida em forma fechada, por uma aproximação linear por partes $\hat{V}_{t+1}(\cdot)$ construída iterativamente. Cada iteração acrescenta um hiperplano de suporte — denominado \emph{corte de Benders} — que preserva a concavidade da aproximação e a torna progressivamente mais precisa \cite{pereira1991}.

\subsection{Etapas \textit{Forward} e \textit{Backward}}
\label{subsec:forward_backward_rev}

A PDDE opera em iterações que alternam duas etapas:

\begin{itemize}
  \item \textbf{Etapa \textit{forward}:} a política corrente é simulada em trajetórias de cenários. Em cada estágio, resolve-se o subproblema com a aproximação disponível de $V_{t+1}$, coletando os estados visitados e produzindo uma estimativa do valor esperado da política (limite superior estocástico).
  \item \textbf{Etapa \textit{backward}:} os estágios são percorridos em ordem reversa. Para cada estado visitado na etapa \textit{forward}, resolve-se o subproblema e utilizam-se as soluções duais para gerar novos cortes de Benders, refinando $\hat{V}_t(\cdot)$ e elevando o limite inferior do valor ótimo.
\end{itemize}

O algoritmo termina quando a diferença entre o limite superior e o limite inferior é inferior a uma tolerância predefinida \cite{pereira1991,shapiro2009}.

\subsection{Vantagens e Limitações}
\label{subsec:vantagens_pdde}

A principal vantagem computacional da PDDE é que o tamanho de cada subproblema de estágio é independente do número total de estágios e cenários, crescendo apenas com a dimensão do vetor de estado. Isso contrasta com a DEQ, cujo tamanho cresce exponencialmente com o horizonte \cite{pereira1991}. Contudo, a complexidade dos subproblemas aumenta com o número de variáveis de estado, o que pode representar uma limitação em problemas com espaços de estado de alta dimensão \cite{shapiro2009}.

O desenvolvimento do pacote \texttt{SDDP.jl} \cite{dowson2021} facilitou a aplicação do método a novos domínios, abstraindo os detalhes de implementação dos cortes de Benders e da estrutura de grafo de política.
